\begin{textsample}{POS Dim 2 – 2011\_c5 – Score -1.00 – 2011\_c5/enade\_2011\_eng\_g\_4\_974c613aa2.txt}  \label{ex:f2_pos_346}
CAPÍTULO 5 DISTRIBUIÇÃO DOS CONCEITOS 5.1 PANORAMA NACIONAL DA DISTRIBUIÇÃO DOS CONCEITOS A Tabela 5.1 apresenta a quantidade e distribuição de cursos de Engenharia – Grupo IV participantes do ENADE/2011, por faixa de conceito e Grande Região.
A diferença entre os cursos tabulados neste capítulo e no capítulo 2 são os cursos sem conceito, em princípio, aqueles sem alunos concluintes que participassem da prova.
Observando -se os dados da Tabela 5.1, nota -se que, dos 117 cursos participantes, 45 ( 38,4 \% ) classificaram -se com conceito 3, o valor modal.
Este foi também o conceito modal nas regiões Norte ( 60,0 \% ), Nordeste ( 38,9 \% ), Sudeste ( 36,7 \% ) e Sul ( 38,5 \% ).
A região Centro-Oeste apresentou a menor moda, conceito 1, com 50,0 \% dos cursos participantes na região.
O conceito 4 foi o segundo mais frequente em nível nacional ( 23,1 \%, correspondendo a 27 cursos ) e o conceito 2, o terceiro ( 21,4 \%, correspondendo a 25 cursos ).
Houve, ainda, 11 cursos ( 9,4 \% ) que receberam conceito 5 e outros sete cursos ( 6,0 \% ) que receberam conceito 1.
Dos 117 cursos de Engenharia – Grupo IV dois ficaram sem conceito ( SC ).
A região Norte participou com cinco cursos ou 4,3 \% do total nacional.
Nesta região, o conceito 3 foi atribuído a três cursos dos cinco participantes, o que equivale a 60,0 \% do total regional.
Nenhum curso ficou sem conceito.
Os demais foram avaliados com conceito 2 ( um curso, correspondendo a 20,0 \% ) e conceito 4 ( também um curso ).
Nenhum curso foi recebeu os conceitos 1 ou 5.
A região Nordeste participou com 18 cursos ou 15,4 \% do total nacional.
Como já comentado, destes, sete cursos, 38,9 \% em termos regionais, obtiveram conceito 3, o conceito modal para a região.
Os conceitos 2 e 4 foram atribuídos a, respectivamente, três e seis cursos ( 16,7 \% e 33,3 \% ).
O conceito 5 foi atribuído a dois cursos ( 11,1 \% ).
Nessa região nenhum dos cursos ficou sem conceito ou recebeu o conceito 1.
Dos 49 cursos participantes da região Sudeste, 18 ( 36,7 \% ) obtiveram conceito 3, o conceito modal.
O conceito 1 foi atribuído a quatro cursos ( 8,2 \% ) e o conceito 2, a 11 ( 22,4 \% ).
Nove cursos ( 18,4 \% ) receberam o conceito 4 e cinco outros ( 10,2 \% ) receberam o conceito 5.
Dois cursos ficaram sem conceito.
Na região Sul a predominância do conceito 3 foi de 38,5 \%, correspondentes a 15 dos 39 cursos participantes na região.
O conceito 2 foi atribuído a dez cursos e o conceito 4, a outros dez cursos ( 25,6 \% para cada conceito ).
O conceito 1 não foi recebido por nenhum curso e o conceito 5, por quatro deles ( 10,3 \% ).
Nenhum dos cursos da região Sul ficou sem conceito.
A metade dos seis cursos participantes na região Centro-Oeste recebeu conceito 1, o conceito modal.
Os demais cursos foram avaliados com conceito 3 ( dois cursos, 33,3 \% ) e conceito 4 ( um curso, 16,7 \% ).
Nesta região nenhum curso ficou sem conceito e tampouco foi avaliado com conceitos 2 ou 5. 5.2 CONCEITOS POR CATEGORIA ADMINISTRATIVA E POR GRANDE REGIÃO A Tabela 5.2 apresenta a distribuição dos cursos participantes do ENADE/2011 de Engenharia – Grupo IV, por Categoria Administrativa, de acordo com os conceitos por eles alcançados, segundo as Grandes Regiões brasileiras.
Dos 117 cursos participantes, 68 ( 58,1 \% ) eram ministrados em instituições públicas e 49 ( 41,9 \% ), em privadas.
De acordo com as informações da Tabela 5.2, em termos nacionais dois cursos ficaram sem conceito um de instituição pública e outro de instituição privadas.
Dos 68 cursos participantes de IES públicas, o conceito 3 foi o valor modal, atribuído a 27 cursos.
Entre os demais cursos participantes, três obtiveram conceito 1 ( 4,4 \% da categoria ), seis receberam conceito 2 ( 8,8 \% ), 23 foram avaliados com conceito 4 ( 33,8 \% ) e oito alcançaram o conceito 5 ( 11,8 \% ).
Na rede privada, o conceito modal foi 2, com 19 cursos dos 49 da categoria.
Entre os demais cursos participantes, quatro receberam conceito 1 e 18 conceito 3.
O conceito 4 foi atribuído a quatro cursos e o conceito 5 a três cursos.
Nesta categoria administrativa, como comentado anteriormente, um dos cursos ficou sem conceito.
Na análise por região, observa -se que na região Norte as instituições privadas participaram com um curso ( 20,0 \% do total regional ), avaliado com conceito 2.
As instituições públicas participaram com quatro cursos ( 80,0 \% do total regional ), dos quais nenhum ficou sem conceito.
O conceito modal para esta Categoria Administrativa, na região Norte, foi o conceito 3, com três cursos, correspondendo a 75,0 \%.
Um curso recebeu o conceito 4, correspondendo a 25,0 \%.
Nenhum curso foi nesta categoria na região recebeu os conceitos 1, 2 ou 5.
Na região Nordeste, a rede privada participou com dois dos 18 cursos, 11,1 \% do total da região, a menor proporção de cursos oferecidos pela rede privada entre todas as regiões.
Nesta Categoria Administrativa os dois únicos cursos receberam o conceito 2.
As instituições públicas dessa região participaram com 16 cursos ( 88,9 \% ), dos quais sete obtiveram o conceito modal, 3.
Os demais cursos foram avaliados com os conceitos : 2 ( um curso ), 4 ( seis cursos ) e 5 ( dois cursos ).
Nenhum curso ficou sem conceito nesta combinação de região e categoria ou recebeu o conceito 1.
Na região Sudeste, a proporção de cursos da rede privada, 55,1 \%, foi mais a mais elevada entre as regiões brasileiras, correspondendo a 27 dos 49 cursos participantes.
Nesta categoria, na região Sudeste, o conceito modal foi 3 ( 11 cursos ).
Os demais foram avaliados com conceito 1 ( dois cursos ), conceito 2 ( nove cursos ), conceito 4 ( dois cursos ), e conceito 5 ( outros dois cursos ).
Nesta combinação de Categoria Administrativa e Grande Região, um curso não recebeu conceito.
Entre os 22 cursos em instituições públicas na região Sudeste, a moda concentrou -se nos conceitos 3 e 4, com sete cursos cada.
Os conceitos 1 e 2 receberam dois cursos cada e três cursos alcançaram o conceito 5.
Na região Sudeste nesta categoria, apenas um curso ficou sem conceito.
As instituições privadas participaram com 17 dos 39 cursos da região Sul, 43,6 \% do total regional.
Nessa região a moda se concentrou nos conceitos 2 e 3 com sete cursos cada.
Os demais foram avaliados com conceito 4 ( dois cursos ) e conceito 5 ( um curso ).
Nesta combinação de Categoria Administrativa e Grande Região, nenhum curso ficou sem receber conceito ou recebeu conceito 1.
As instituições públicas na região Sul participaram com 22 cursos ( 56,4 \% ), dos quais nenhum ficou sem conceito.
Os conceitos modais foram 3 e 4, com oito cursos cada.
Além disso, os conceitos 2 e 5 foram recebidos por três cursos cada.
Na região Centro-Oeste, dois dos seis cursos participantes eram de instituições privadas ( 33,3 \% em termos regionais ).
Ambos receberam conceito 1, conceito modal.
Dos quatro cursos de instituições públicas, dois foram avaliados no conceito 3, conceito modal e os outros dois receberam os conceitos 1 e 4.
Nesta região, nenhum curso de IES pública ficou sem conceito ou recebeu os conceitos 2 ou 5. 5.3 CONCEITOS POR ORGANIZAÇÃO ACADÊMICA E POR GRANDE REGIÃO Na Tabela 5.3 encontra -se a distribuição dos conceitos atribuídos aos cursos participantes do ENADE/2011 na Área de Engenharia – Grupo IV, por Organização Acadêmica, segundo as Grandes Regiões brasileiras.
Dos 117 cursos participantes, 91 eram oferecidos em Universidades, 10 em Centros Universitários e os demais 16 em Faculdades.
Esta distribuição corresponde a, respectivamente, 77,8 \%, 8,5 \% e 13,7 \% dos cursos.
De acordo com os dados apresentados, oito dos 11 cursos avaliados com conceito 5 eram vinculados a Universidades.
Apenas um curso em Universidades ficou sem conceito e o conceito 3 foi o modal, com 38 cursos.
Os demais cursos avaliados receberam os conceitos : 1 ( cinco cursos ), 2 ( 14 cursos ), 4 ( 25 cursos ) e conceito 5 ( oito cursos, como já mencionado ).
Entre os cursos em Centros Universitários, a moda se concentrou nos conceitos 2 e 3 com três cursos cada.
Neste tipo de organização acadêmica um curso ficou sem conceito e os outros cursos receberam os conceitos 1, 4 e 5, com um curso cada.
Nas Faculdades, nenhum dos 16 cursos ficou sem conceito e oito receberam o conceito modal 2.
Dos demais cursos neste tipo de Organização Acadêmica, um recebeu conceito 1, quatro conceito 3, outro curso o conceito 4 e dois o conceito 5.
Nenhum curso ficou sem conceito.
Considerando -se separadamente as regiões brasileiras, verifica -se que na região Norte as Universidades concentraram quatro dos cinco cursos participantes.
Dos cursos em Universidades, nenhum ficou sem conceito ou recebeu conceitos 1, 2 ou 5.
Três cursos receberam o conceito 3 e o curso restante recebeu o conceito 4.
Os Centros Universitários da região Norte participaram com um curso que recebeu o conceito 2.
As Faculdades não tiveram nenhum curso participante nesta região.
Na região Nordeste, todos os 18 cursos na Área de Engenharia – Grupo IV eram de Universidades.
Nenhum deles ficou sem conceito e o conceito modal foi 3, com sete cursos.
Os demais receberam conceito 2 ( três cursos ), 4 ( seis cursos ) e conceito 5 ( dois cursos ).
Na região Sudeste, as Universidades concentraram 28 dos 49 cursos da região.
Entre os cursos oferecidos em Universidades na região, o conceito modal foi 3 com 11 cursos e um curso ficou sem conceito.
Os conceitos 2 e 4 receberam, respectivamente, quatro e sete cursos.
Os conceitos extremos, 1 e 5, receberam três e dois cursos respectivamente.
Os Centros Universitários participaram com sete cursos na região Sudeste, dos quais três obtiveram conceito modal, 3, e um ficou sem conceito.
Os demais receberam os conceitos 2, 3 e 4 com um curso cada.
As Faculdades foram representadas por 14 cursos na região Sudeste, que se distribuíram nos conceitos : 1 ( um curso ), 2 ( seis cursos, conceito modal ), 3 ( quatro cursos ), 4 ( um curso ) e 5 ( dois cursos ).
Nenhum curso oferecido em Faculdade ficou sem conceito.
Dos 39 cursos da região Sul, 36 eram oferecidos em Universidades, 92,3 \% dos cursos da região.
O conceito modal foi 3, com 15 cursos.
Nesse tipo de organização, nenhum dos cursos ficou sem conceito ou foi alocado ao conceito 1.
Os demais receberam os conceitos : 2 ( sete cursos ), 4 ( dez cursos ) e 5 ( quatro cursos ).
Os Centros Universitários da região Sul tiveram apenas um curso participante que recebeu conceito 2.
As Faculdades na região Sul participaram com dois cursos e ambos receberam conceito 2.
Na região Centro-Oeste, cinco dos seis cursos eram de Universidades.
Nesse tipo de organização, nenhum curso ficou sem conceito.
A moda se concentrou nos conceitos 1 e 3 com dois cursos cada e um curso foi alocado ao conceito 4.
Os Centros Universitários da região Centro-Oeste contaram com um curso, que recebeu o conceito 1.
Não houve participação de curso oferecido em Faculdades na região Centro-Oeste.

% matched lemmas: 
\end{textsample}
